% Auto-generated from book/07_representations_z_flux.md — do not edit by hand
% Regenerated by scripts/md_to_latex.py

\chapter{The \(Z \mapsto\) Flux Map and Representation Theory}
\label{ch:ch07_z_map}

This chapter connects the classified flux topographs and Magic Islands of Chapters 5--6 to the atomic number \(Z\). We define the \textbf{\(Z \mapsto\) flux map}, interpret it as a representation-theoretic bridge (which ``numbers'' are represented by stable flux configurations), and examine how Magic Islands align with chemically and physically special elements. The construction remains a \textbf{Model} with explicit validation pathways (Chapter 10).

\textbf{Learning goals}

\begin{enumerate}
\item Define the \(Z \mapsto\) flywheel / flux configuration map formally.  
\item Interpret the map as representation theory on the gauged Hopf lattice.  
\item Connect Magic Islands to the periodic table and noble-gas stability.  
\item Examine chemical and physical proxies (stability scores, shell descriptors, extended fields).  
\item Prepare the bridge to class-group ideas (Chapter 8) and observational validation (Chapter 10).
\end{enumerate}

\textbf{Figures in this chapter}

\begin{center}
\small
\setlength{\tabcolsep}{4pt}
\begin{tabularx}{\textwidth}{@{}>{\raggedright\arraybackslash}p{0.12\textwidth}>{\raggedright\arraybackslash}X>{\raggedright\arraybackslash}X@{}}
\toprule
Tag & File & Role \\
\midrule
Fig. 7.1 & \texttt{figures/\allowbreak{}fig7\_\allowbreak{}1\_\allowbreak{}z\_\allowbreak{}to\_\allowbreak{}flywheel.\allowbreak{}png} & Schematic of the \(Z \mapsto\) flux map \\
Fig. 7.2 & \texttt{figures/\allowbreak{}fig7\_\allowbreak{}2\_\allowbreak{}magic\_\allowbreak{}island\_\allowbreak{}periodic\_\allowbreak{}table.\allowbreak{}png} & Magic Island scores on a schematic table \\
Fig. 7.3 & \texttt{figures/\allowbreak{}fig7\_\allowbreak{}3\_\allowbreak{}stability\_\allowbreak{}vs\_\allowbreak{}z.\allowbreak{}png} & \texttt{stability\_\allowbreak{}score} vs \(Z\) (portal data) \\
Fig. 7.4 & \texttt{figures/\allowbreak{}fig7\_\allowbreak{}4\_\allowbreak{}electron\_\allowbreak{}cloud\_\allowbreak{}flux.\allowbreak{}png} & Electron cloud as flux distribution \\
Aux A7.1 & \texttt{figures/\allowbreak{}aux7\_\allowbreak{}1\_\allowbreak{}he\_\allowbreak{}fe\_\allowbreak{}au\_\allowbreak{}stills.\allowbreak{}png} & He, Fe, Au flywheel stills \\
\bottomrule
\end{tabularx}
\end{center}


\textbf{Claim discipline}

\begin{center}
\small
\setlength{\tabcolsep}{4pt}
\begin{tabularx}{\textwidth}{@{}>{\raggedright\arraybackslash}p{0.70\textwidth}>{\raggedright\arraybackslash}X@{}}
\toprule
Claim & Type \\
\midrule
Four-square theorem and classical representation by norms (Ch. 1) & \textbf{Theorem} \\
Existence and implementation of the coded \(Z \mapsto\) map & \textbf{Software fact} \\
Map as representation theory; Magic Islands as periodic-table proxy; physical emergence & \textbf{Model} \\
Any claim that topology \emph{derives} real chemistry as a law of nature & \textbf{Hypothesis} (Ch. 10 validation) \\
\texttt{map\_\allowbreak{}z\_\allowbreak{}to\_\allowbreak{}flywheel[\_\allowbreak{}extended]} metrics & \textbf{Software fact} \\
\bottomrule
\end{tabularx}
\end{center}


\bigskip
\noindent\rule{\textwidth}{0.4pt}
\bigskip

\section{The \(Z \mapsto\) flux map}
\label{ch:ch07_z_map:the-z-mapsto-flux-map}

The Kingdom Come portal implements a function that takes an atomic number \(Z\) and returns a flywheel configuration together with stability and chemistry-facing metrics:

\begin{align*}
Z \longmapsto \bigl(&
  \text{flywheel / detuning state},\\
  &\texttt{stability\_score},\;
  \texttt{stability\_class},\\
  &\texttt{is\_noble\_gas},\;\ldots
\bigr).
\end{align*}

\textbf{Primary entry points}

\begin{center}
\small
\setlength{\tabcolsep}{4pt}
\begin{tabularx}{\textwidth}{@{}>{\raggedright\arraybackslash}X>{\raggedright\arraybackslash}X@{}}
\toprule
Function & Role \\
\midrule
\texttt{map\_\allowbreak{}z\_\allowbreak{}to\_\allowbreak{}flywheel(z)} & Core map: detuning, gauge params, \texttt{stability\_\allowbreak{}score}, class labels \\
\texttt{map\_\allowbreak{}z\_\allowbreak{}to\_\allowbreak{}flywheel\_\allowbreak{}extended(z)} & Adds alignment, IE proxies, magnetic-moment validation fields \\
\bottomrule
\end{tabularx}
\end{center}


There is still \textbf{no} \texttt{magic\_\allowbreak{}flag} key. Treat high \texttt{stability\_\allowbreak{}score}, noble-gas class strings, and \texttt{is\_\allowbreak{}noble\_\allowbreak{}gas=True} as the Model's island indicators.

The map is built on the gauged Hopf lattice and flux-topograph language of previous chapters. Different detuning parameters select different stability islands; a documented ultra-stable lock (Magic Island Sweep reference in portal notes) serves as a noble-gas-like template.

\begin{figure}[htbp]
  \centering
  \includegraphics[width=0.92\textwidth,height=0.42\textheight,keepaspectratio]{fig7_1_z_to_flywheel.png}
  \caption{Pipeline: atomic number \(\rightarrow\) lattice / gauge data \(\rightarrow\) flux flywheel + topograph \(\rightarrow\) scored metrics.}
  \label{fig:fig7-1-z-to-flywheel}
\end{figure}

\textbf{Claim type.} Reproducible coded map: \textbf{Software fact}. Interpreting outputs as \emph{the} emergence mechanism of the periodic table: \textbf{Model}.

\subsection{Book helper}
\label{ch:ch07_z_map:book-helper}

\begin{lstlisting}[style=qga]
stability_landscape_z(z_values=...)          # list of dict rows
stability_landscape_z(z_range=(1, 50))       # inclusive range
stability_landscape_z(z_range=(1, 30), extended=True)
\end{lstlisting}

\bigskip
\noindent\rule{\textwidth}{0.4pt}
\bigskip

\section{Representation theory on the gauged Hopf lattice}
\label{ch:ch07_z_map:representation-theory-on-the-gauged-hopf-lattice}

In classical number theory, representation by quadratic forms asks which integers \(n\) satisfy \(f(x,y)=n\) for a given form (Hatcher Chapter 6). Quaternion norms answer which \(n\) are sums of four squares (Chapter 1, \textbf{Theorem}).

Here we ask an analogous \textbf{Model} question:

> Which atomic numbers \(Z\) are ``represented'' by stable flux configurations on the gauged Hopf lattice?

Stable flywheels correspond to reduced configurations with high \texttt{magic\_\allowbreak{}island\_\allowbreak{}score} (Chapter 6) and high portal \texttt{stability\_\allowbreak{}score}. The map therefore selects, for each \(Z\), a configuration (or orbit) intended to match observed chemical/nuclear specialness.

\begin{center}
\small
\setlength{\tabcolsep}{4pt}
\begin{tabularx}{\textwidth}{@{}>{\raggedright\arraybackslash}X>{\raggedright\arraybackslash}X@{}}
\toprule
Classical & Quaternionic / Hopf Model \\
\midrule
Form \(f\) & Lattice + gauge group + flux functional \\
Represented integer \(n\) & Atomic number \(Z\) (and related quantum numbers) \\
Class / genus of forms & Equivalence class of topographs / flywheels \\
Class number & \texttt{class\_\allowbreak{}number\_\allowbreak{}analogue} / island counts (Ch. 6) \\
\bottomrule
\end{tabularx}
\end{center}


This is a \textbf{representation-theoretic reading}, not a theorem that every chemical property is a lattice invariant.

\begin{figure}[htbp]
  \centering
  \includegraphics[width=0.92\textwidth,height=0.42\textheight,keepaspectratio]{fig7_2_magic_island_periodic_table.png}
  \caption{Cells colored by portal \texttt{stability\_\allowbreak{}score}; blue borders mark noble-gas \(Z\). Layout is a simplified table for visualization (f-block truncated)---software + Model.}
  \label{fig:fig7-2-magic-island-periodic-table}
\end{figure}

\bigskip
\noindent\rule{\textwidth}{0.4pt}
\bigskip

\section{Chemistry-facing scores (Hypothesis --- developed in Part V)}
\label{ch:ch07_z_map:chemistry-facing-scores-hypothesis-developed-in-part-v}

Magic Islands as \emph{chemical} stability, noble-gas locks, and ionization-energy alignment are \textbf{not} theorems of Part III. They are Hypothesis H2 / H3, with Table T4 protocols and an H2 ablation (noble-gas bonuses off) in \textbf{Chapter 10} and Appendix D. This section keeps only a pointer so the \(Z\mapsto\) map can be used as a \textbf{Model} representation without promoting chemistry.

Figure 7.3 is a portal software curve, not a finding.

\begin{figure}[htbp]
  \centering
  \includegraphics[width=0.92\textwidth,height=0.42\textheight,keepaspectratio]{fig7_3_stability_vs_z.png}
  \caption{Portal \texttt{stability\_\allowbreak{}score} vs atomic number. Gold: noble-gas \(Z\) (explicit bonuses live in the portal map). Interpretation: \textbf{Hypothesis}, Part V. Lab 7 must ablate those bonuses before any alignment claim.}
  \label{fig:fig7-3-stability-vs-z}
\end{figure}

\textbf{Open Problem 4 (shared with Ch. 10).} Up to gauge, is the map \(Z\mapsto\) flywheel configuration unique under stated axioms? See \texttt{notes/\allowbreak{}open\_\allowbreak{}problems.\allowbreak{}md}.

\bigskip
\noindent\rule{\textwidth}{0.4pt}
\bigskip

\section{Electron-cloud visuals (Software fact --- Part V)}
\label{ch:ch07_z_map:electron-cloud-visuals-software-fact-part-v}

Portal electron-cloud pictures (\texttt{kingdom.\allowbreak{}viz.\allowbreak{}electron\_\allowbreak{}cloud}) are \textbf{Software facts} about the renderer, not Schrödinger solutions and not a Part III theorem. Chemistry-facing reading: \textbf{Part V}.

\begin{figure}[htbp]
  \centering
  \includegraphics[width=0.92\textwidth,height=0.42\textheight,keepaspectratio]{fig7_4_electron_cloud_flux.png}
  \caption{Schematic closed-shell vs mid-table flux-derived profiles (portal style).}
  \label{fig:fig7-4-electron-cloud-flux}
\end{figure}

\begin{figure}[htbp]
  \centering
  \includegraphics[width=0.92\textwidth,height=0.42\textheight,keepaspectratio]{aux7_1_he_fe_au_stills.png}
  \caption{Portal Flux Flywheel stills: He (\(Z=2\)), Fe (\(Z=26\)), Au (\(Z=79\)). Sources: \texttt{kingdom/\allowbreak{}z\_\allowbreak{}knowns/\allowbreak{}frame\_\allowbreak{}0002.\allowbreak{}png}, \texttt{frame\_\allowbreak{}0026.\allowbreak{}png}, \texttt{frame\_\allowbreak{}0079.\allowbreak{}png}.}
  \label{fig:aux7-1-he-fe-au-stills}
\end{figure}

\bigskip
\noindent\rule{\textwidth}{0.4pt}
\bigskip

\section{First computational labs}
\label{ch:ch07_z_map:first-computational-labs}

Portal: \texttt{map\_\allowbreak{}z\_\allowbreak{}to\_\allowbreak{}flywheel[\_\allowbreak{}extended]} \(\cdot\) \texttt{stability\_\allowbreak{}landscape\_\allowbreak{}z} \(\cdot\) \textbf{Appendix C \S{}C.3}.

\begin{itemize}
\item \textbf{7.A} Scores for \(Z\in\{2,10,26,79,118\}\).
\item \textbf{7.B} \texttt{stability\_\allowbreak{}landscape\_\allowbreak{}z(z\_\allowbreak{}range=(1,50))} high islands.
\item \textbf{7.C} Noble vs neighbors (\(Z=10,11,12\)).
\item \textbf{7.D} Gradio Flux Flywheel / Aux A7.1.
\item \textbf{7.E} Extended fields He vs Fe.
\end{itemize}

\begin{lstlisting}[style=qga,language=python]
from kingdom.core.flux_flywheel import map_z_to_flywheel
print(map_z_to_flywheel(2)["stability_score"], map_z_to_flywheel(2).get("is_noble_gas"))
\end{lstlisting}

\bigskip
\noindent\rule{\textwidth}{0.4pt}
\bigskip

\section{Exercises}
\label{ch:ch07_z_map:exercises}

\textbf{7.A (hand).} In one paragraph, explain why the \(Z\mapsto\) map can be viewed as a representation problem on the gauged Hopf lattice.

\textbf{7.B (hand).} Why might Magic Islands near noble-gas \(Z\) be especially stable in the Model? Connect to closed-shell ideas and Chapter 6 reduced configurations.

\textbf{7.C (code).} Complete Labs 7.A--7.B. Tabulate \texttt{stability\_\allowbreak{}score} and \texttt{is\_\allowbreak{}noble\_\allowbreak{}gas} for \(Z=1\) to \(30\). Identify the strongest islands.

\textbf{7.D (code).} Complete Lab 7.E for \(Z=2\) and \(Z=26\). What differences appear in extended fields?

\textbf{7.E (visual).} Render or inspect flux-derived clouds for He and Fe. Describe how Model geometry differs between a closed-shell case and a mid-table case.

\textbf{7.F (Hatcher bridge).} In Hatcher Chapter 6, representation by quadratic forms uses genus, characters, and quadratic reciprocity. Sketch how the \(Z\mapsto\) map might eventually use ``genus'' of flux configurations or topological characters on the lattice.

\textbf{7.G (forward).} Why will class-group ideas (Chapter 8) be relevant to which \(Z\) land in Magic Islands?

\textbf{7.H (software honesty).} Distinguish in one paragraph: (i) the coded \(Z\mapsto\) map as software, (ii) its use as a periodic-table proxy (\textbf{Model}), (iii) any claim that it derives real chemistry from topology (\textbf{Hypothesis}).

\textbf{7.I (Open Problem 4 teaser).} Fix a small set of gauge sequences from Chapter 4. For two nearby \(Z\), inspect whether flywheel metrics could plausibly share an orbit. What would count as non-uniqueness of the map?

\bigskip
\noindent\rule{\textwidth}{0.4pt}
\bigskip

\section{Code and asset pointers}
\label{ch:ch07_z_map:code-and-asset-pointers}

\begin{lstlisting}[style=qga]
kingdom.core.flux_flywheel.map_z_to_flywheel[_extended]
qga/lib/flux_topograph.stability_landscape_z
kingdom.viz.magic_island.build_magic_island_heatmap
kingdom.viz.electron_cloud.build_electron_cloud_figure
kingdom/z_knowns/frame_XXXX.png   # flywheel stills
\end{lstlisting}

\textbf{Figures:} \texttt{scripts/\allowbreak{}generate\_\allowbreak{}ch7\_\allowbreak{}figures.\allowbreak{}py} \textbf{Related portal:} Flux Flywheel tab (interactive \(Z\) slider). \textbf{Open problems:} OP4 (uniqueness); OP3 (island invariants); OP5 (\(350/\pi\), Ch. 10).

\bigskip
\noindent\rule{\textwidth}{0.4pt}
\bigskip

\section{Looking ahead}
\label{ch:ch07_z_map:looking-ahead}

The \(Z\mapsto\) flux map links the classified arithmetic objects of Chapters 5--6 to concrete atomic numbers. In \textbf{Chapter 8} we lift the class group itself to the quaternionic / Hopf setting, giving a deeper arithmetic home for Magic Island counts and equivalence classes. \textbf{Chapter 10} returns to observational validation of the whole construction---including statistical protocols for the \(350/\pi\) signature and \(Z\)-map correlations.

With the representation bridge in place, we are ready for the ideal-class-group lift.

\bigskip
\noindent\rule{\textwidth}{0.4pt}
\bigskip

